% TEMPLATE-MILC-v3.tex - plantilla maestra para todas las guias MILC v3
% Ver scripts/generadores/build-guia-milc.py para la lista completa de
% claves esperadas. El script valida que no queden placeholders sin
% reemplazar antes de escribir el archivo final.

\documentclass[11pt,letterpaper]{article}
\usepackage[margin=1.75cm]{geometry}
\usepackage[table]{xcolor}
\usepackage{fontspec}
\usepackage[spanish]{babel}
\usepackage{tikz}
\usepackage[most]{tcolorbox}
\usepackage{tabularx}
\usepackage{array}
\usepackage{fancyhdr}
\usepackage{titlesec}
\usepackage{ragged2e}
\usepackage{enumitem}
\usepackage{microtype}

\setmainfont{Helvetica Neue}
\setsansfont{Helvetica Neue}
\setlength{\parindent}{0pt}
\setlength{\parskip}{6pt}
\hyphenpenalty=200
\exhyphenpenalty=200
\pretolerance=400
\tolerance=2000
\setlength{\emergencystretch}{1em}
\renewcommand{\arraystretch}{1.25}
\setlist{leftmargin=5mm,itemsep=1mm,topsep=1mm}
\newcolumntype{Y}{>{\RaggedRight\arraybackslash}X}

\definecolor{milcMagenta}{HTML}{E6007E}
\definecolor{milcVino}{HTML}{5A0038}
\definecolor{milcTurquesa}{HTML}{0093A5}
\definecolor{milcVerde}{HTML}{58A923}
\definecolor{milcMostaza}{HTML}{E5B400}
\definecolor{milcOcre}{HTML}{B86600}
\definecolor{milcNegro}{HTML}{241816}
\definecolor{milcGris}{HTML}{F7F7F4}
\definecolor{milcLinea}{HTML}{E8E2DD}
\definecolor{milcGradePrimary}{HTML}{84CC16}
\definecolor{milcGradeDark}{HTML}{4D7C0F}
\definecolor{milcGradeSoft}{HTML}{ECFCCB}
\definecolor{milcGradeAccent}{HTML}{CFE7FF}
\definecolor{milcGradeText}{HTML}{FFFFFF}
\color{milcNegro}

\pagestyle{fancy}
\fancyhf{}
\lhead{\small Tecnología e Informática · Grado 7}
\rhead{\small Guía 4 · MILC v3}
\cfoot{\small\thepage}
\renewcommand{\headrulewidth}{0pt}

\newtcolorbox{titlebox}[1]{
  enhanced,colback=#1,colframe=#1,arc=7mm,boxrule=0pt,
  left=7mm,right=7mm,top=3mm,bottom=3mm,
  before skip=8pt,after skip=8pt,
  fontupper=\sffamily\bfseries\Large\color{white}
}

\newtcolorbox{softbox}[3]{
  enhanced,colback=#2,colframe=#1,borderline west={4pt}{0pt}{#1},
  arc=2mm,boxrule=.4pt,left=4mm,right=4mm,top=3mm,bottom=3mm,
  before skip=6pt,after skip=8pt,title={#3},coltitle=white,
  colbacktitle=#1,fonttitle=\sffamily\bfseries,fontupper=\RaggedRight,
  attach boxed title to top left={xshift=3mm,yshift=-2mm},
  boxed title style={arc=2mm, boxrule=0pt}
}

\newtcolorbox{drawbox}[2]{
  enhanced,colback=white,colframe=#1,arc=4mm,boxrule=1pt,
  left=5mm,right=5mm,top=4mm,bottom=4mm,before skip=8pt,after skip=8pt,
  title={#2},coltitle=white,colbacktitle=#1,fonttitle=\sffamily\bfseries,
  fontupper=\RaggedRight,
  attach boxed title to top left={xshift=4mm,yshift=-2mm},
  boxed title style={arc=3mm, boxrule=0pt}
}

\newtcolorbox{infoband}[1]{
  enhanced,colback=#1!8,colframe=#1!50,arc=2mm,boxrule=.5pt,
  left=4mm,right=4mm,top=2mm,bottom=2mm,before skip=4pt,after skip=4pt,
  fontupper=\small\RaggedRight
}

% Puente narrativo entre secciones (contrato editorial Conéctate).
% Da continuidad: "Acabas de X. Ahora vas a Y porque Z."
% Visualmente: banda izquierda en color del grado, texto en itálica pequeña.
\newtcolorbox{puentebox}{
  enhanced,colback=milcGradeSoft,colframe=milcGradeDark,
  borderline west={3pt}{0pt}{milcGradeDark},
  arc=0mm,boxrule=0pt,left=5mm,right=4mm,top=2.5mm,bottom=2.5mm,
  before skip=4pt,after skip=8pt,
  fontupper=\itshape\small\color{milcNegro}\RaggedRight
}
\newcommand{\puente}[1]{%
  \begin{puentebox}%
    \textcolor{milcGradeDark}{\textbf{\textrightarrow}}\hspace{2mm}#1%
  \end{puentebox}%
}

\newcommand{\linead}[1]{\noindent\color{milcLinea}\rule{#1}{0.5pt}\color{milcNegro}}
\newcommand{\respuesta}[1]{\par\vspace{2mm}\foreach \n in {1,...,#1}{\noindent\color{milcLinea}\rule{\linewidth}{0.5pt}\par\vspace{5mm}}\color{milcNegro}}
\newcommand{\checkbox}{\textcolor{milcGradeDark}{\fbox{\phantom{x}}}\hspace{5pt}}

\newcommand{\phasecard}[4]{
\begin{tcolorbox}[enhanced,colback=#3,colframe=#2,arc=3mm,boxrule=.5pt,height=3.05cm,valign=center,halign=center,left=1.5mm,right=1.5mm,top=2mm,bottom=2mm]
{\sffamily\bfseries\fontsize{22}{24}\selectfont\textcolor{#2}{#1}}\par
{\sffamily\bfseries\small #4}
\end{tcolorbox}
}

\begin{document}
\thispagestyle{empty}
\begin{tikzpicture}[remember picture,overlay]
  \fill[white] (current page.south west) rectangle (current page.north east);
  \path[fill=milcGradePrimary,rounded corners=16mm]
    ([xshift=.55cm,yshift=-.55cm]current page.north west) rectangle
    ([xshift=-.55cm,yshift=3.65cm]current page.south east);

  \node[anchor=north west,text=milcGradeText] at ([xshift=2.0cm,yshift=-1.85cm]current page.north west)
    {\sffamily\bfseries\fontsize{14}{16}\selectfont GRADO};
  \node[anchor=north west,text=milcGradeText,opacity=.82] at ([xshift=2.0cm,yshift=-2.75cm]current page.north west)
    {\sffamily\bfseries\fontsize{9}{11}\selectfont SERIE GUÍAS · TECNOLOGÍA E INFORMÁTICA};

  \begin{scope}[shift={([xshift=-3.05cm,yshift=-2.55cm]current page.north east)}]
    \fill[white,opacity=.80] (-.62,-.52) rectangle (.62,.52);
    \draw[milcGradeText,line width=1pt] (-.62,-.52) rectangle (.62,.52);
    \fill[milcGradeDark] (-.42,-.38) rectangle (-.22,.06);
    \fill[milcGradeAccent] (-.10,-.38) rectangle (.10,.24);
    \fill[milcGradePrimary!60!black] (.22,-.38) rectangle (.42,.38);
  \end{scope}

  \node[anchor=west,text=milcGradeText] at ([xshift=1.9cm,yshift=-12.85cm]current page.north west)
    {\sffamily\bfseries\fontsize{124}{124}\selectfont 7};
  \draw[milcGradeText,line width=1.7pt]
    ([xshift=4.52cm,yshift=-11.68cm]current page.north west) circle (.13cm);

  \node[anchor=west,text=milcGradeText,text width=8.7cm,align=left] at ([xshift=10.35cm,yshift=-11.6cm]current page.north west)
    {
     {\sffamily\bfseries\fontsize{19}{22}\selectfont Word, Excel\\y PowerPoint\\en línea ---\\las 3 estrellas\\del paquete}\\[4mm]
     {\sffamily\bfseries\fontsize{23}{26}\selectfont Guía 4-7-TIC}\\[2mm]
     {\sffamily\fontsize{13}{16}\selectfont Período 1 · Trabajo colaborativo en la nube}\\[5mm]
     {\sffamily\bfseries\fontsize{11}{13}\selectfont Institución Educativa Sor María Juliana}\\[1mm]
     {\sffamily\fontsize{10}{12}\selectfont Autor: Dr. Álvaro Cárdenas Orozco}
    };

  \path[fill=milcGradeSoft,rounded corners=7mm]
    ([xshift=1.35cm,yshift=.85cm]current page.south west) rectangle
    ([xshift=-1.35cm,yshift=4.25cm]current page.south east);
  \draw[milcGradeDark,line width=.65pt,rounded corners=7mm]
    ([xshift=1.35cm,yshift=.85cm]current page.south west) rectangle
    ([xshift=-1.35cm,yshift=4.25cm]current page.south east);
  \node[anchor=center,text=milcNegro,text width=17.0cm,align=center] at ([yshift=2.55cm]current page.south)
    {\sffamily\fontsize{10}{12}\selectfont Metodología MILC · Apertura ancestral · Pensamiento computacional · Triángulo Dussel-Estoicismo-Floridi\\
    \textcolor{milcGradeDark}{\bfseries Producto:} Tres archivos nuevos en tu OneDrive: \textbf{(a) Word en línea} con un texto de 1 página sobre la minga digital (continuación del tema de la S1); \textbf{(b) Excel en línea} con una tabla de 5 columnas y 10 filas de calificaciones inventadas + fórmula \texttt{=PROMEDIO()} aplicada; \textbf{(c) PowerPoint en línea} con 5 diapositivas presentando \emph{``Mis 3 apps favoritas de Microsoft 365''}. Cada archivo guardado en \texttt{Mis-documentos/Colegio/Tecnologia/Periodo-1}. Más \textbf{en el cuaderno}: tabla comparativa de las 3 apps (Word/Excel/PowerPoint) con 5 criterios.};
\end{tikzpicture}
\mbox{}
\newpage

\section*{Apertura: Word, Excel y PowerPoint en línea --- las 3 herramientas estrella}

\begin{softbox}{milcGradeDark}{milcGradeSoft}{Saber ancestral}
\textbf{En el taller de don Lucho el relojero de Cartago había 3 herramientas estrella que él usaba el 80\% del tiempo.} A pesar de que tenía en su vitrina de la calle 14 \emph{decenas de herramientas} para reparar relojes, en la práctica cotidiana siempre echaba mano a 3: \textbf{la lupa de aumento} para ver lo pequeño, \textbf{las pinzas delicadas} para agarrar piezas diminutas, y \textbf{el destornillador miniatura} para abrir tapas. \emph{Las otras herramientas (martillo, aceite, pulidor) las usaba solo en casos especiales}. Cuando un aprendiz nuevo llegaba, don Lucho le decía: \emph{``Primero domine estas 3. Después se da cuenta de que con ellas hace el 80\% del oficio. Las otras se aprenden con el tiempo.''} \textbf{Hoy en tu taller digital pasa lo mismo}: aunque Microsoft 365 tenga 8 aplicaciones, en la práctica usas 3 el 80\% del tiempo: \emph{Word, Excel, PowerPoint}. Dominar bien esas 3 te hace eficiente. Después aprenderás las otras 5.
\end{softbox}

\begin{softbox}{milcTurquesa}{milcGris}{Saber contemporáneo}
\textbf{Word, Excel y PowerPoint son las 3 aplicaciones más usadas} del paquete Microsoft 365. Cada una tiene una versión \textbf{en línea} (\emph{Word Online}, \emph{Excel Online}, \emph{PowerPoint Online}) que funciona \textbf{dentro del navegador} (sin instalar nada en tu PC). \textbf{Las versiones en línea tienen 95\% de las funciones} que las versiones de escritorio. Para 7° grado, lo en línea sobra. \textbf{Ventaja clave de la versión en línea:} se guarda automáticamente en OneDrive \emph{cada pocos segundos}. Si se cae internet o se apaga el computador, no pierdes nada. Eso resuelve el problema de \emph{``olvidé guardar y se apagó''} que vivían tus papás. \textbf{Las 3 apps cubren 3 propósitos distintos:} Word para \emph{textos}, Excel para \emph{números}, PowerPoint para \emph{presentar}. Hoy creas uno de cada uno para tomarles confianza.
\end{softbox}

\begin{softbox}{milcMostaza}{milcGris}{Pregunta-puente}
\textit{Si te dieran a escoger \textbf{una sola app} de las 3 para aprender muy bien (a costa de no usar las otras 2), ¿\textbf{cuál escogerías} y por qué? ¿Word, Excel o PowerPoint? La respuesta dice mucho sobre el oficio que más te llama.}
\end{softbox}

\begin{softbox}{milcVerde}{milcGris}{Saber y saber hacer}
Al terminar la clase: (1) podrás \textbf{identificar} las diferencias clave entre Word, Excel y PowerPoint; (2) sabrás \textbf{aplicar} cada una a su tipo de tarea; (3) podrás \textbf{evaluar} cuál usar para una tarea nueva; (4) habrás \textbf{creado} 1 archivo de cada uno en tu OneDrive.
\end{softbox}

\small
\begin{tabularx}{\textwidth}{>{\bfseries}p{3.0cm}Y}
\rowcolor{milcGradePrimary!12}Autor & Dr. Álvaro Cárdenas Orozco\\
DBA & Utiliza herramientas digitales para trabajar de manera colaborativa con otros (MEN, T\&I 7°)\\
Referentes & Saber ancestral de la minga · Microsoft 365 y OneDrive · Coautoría asincrónica\\
Producto final & Tres archivos nuevos en tu OneDrive: \textbf{(a) Word en línea} con un texto de 1 página sobre la minga digital (continuación del tema de la S1); \textbf{(b) Excel en línea} con una tabla de 5 columnas y 10 filas de calificaciones inventadas + fórmula \texttt{=PROMEDIO()} aplicada; \textbf{(c) PowerPoint en línea} con 5 diapositivas presentando \emph{``Mis 3 apps favoritas de Microsoft 365''}. Cada archivo guardado en \texttt{Mis-documentos/Colegio/Tecnologia/Periodo-1}. Más \textbf{en el cuaderno}: tabla comparativa de las 3 apps (Word/Excel/PowerPoint) con 5 criterios.\\
\end{tabularx}
\normalsize

\puente{Hoy haces 4 cosas: distingues Word vs Excel vs PowerPoint, abres cada una en línea, creas un archivo en cada una, guardas todo en OneDrive.}

\begin{titlebox}{milcGradeDark}
Ruta MILC: así vamos a aprender
\end{titlebox}
\begin{tabularx}{\textwidth}{>{\centering\arraybackslash}X>{\centering\arraybackslash}X>{\centering\arraybackslash}X>{\centering\arraybackslash}X}
\phasecard{1}{milcVerde}{milcGris}{Escuta\\[-1mm]\small Observo, escucho y pregunto.} &
\phasecard{2}{milcTurquesa}{milcGris}{Sistematización\\[-1mm]\small Pensamiento computacional.} &
\phasecard{3}{milcMagenta}{milcGris}{Praxis\\[-1mm]\small Construyo y aplico.} &
\phasecard{4}{milcVino}{milcGris}{Evaluación\\[-1mm]\small Reflexión liberadora.}
\end{tabularx}


\puente{Empezamos con un test de 9 situaciones: ¿Word, Excel o PowerPoint?}

\begin{titlebox}{milcVerde}
Fase 1 · Escuta
\end{titlebox}

\begin{softbox}{milcVerde}{milcGris}{Escena para escuchar}
\textbf{Actividad 1 · IDENTIFICA --- 9 situaciones, 3 apps} (10 min · individual). En el cuaderno, escribe esta lista del 1 al 9. Para cada situación, escribe \textbf{W (Word)}, \textbf{E (Excel)} o \textbf{P (PowerPoint)} según cuál usarías. \emph{(1) Escribir un cuento de 2 páginas. (2) Llevar las notas de los 30 compañeros de la clase con promedios. (3) Hacer una presentación oral con 12 diapositivas sobre los animales del Valle. (4) Redactar una carta formal al rector. (5) Hacer un cuadro de gastos del mes. (6) Crear una invitación a un cumpleaños con diseño bonito. (7) Escribir un poema sobre Cartago. (8) Llevar un control de las horas que estudias cada día. (9) Mostrar a tu clase los resultados de una encuesta con gráficos.}
\end{softbox}

\checkbox Listé las 9 situaciones numeradas.\\
\checkbox Marqué W, E o P al lado de cada una.\\
\checkbox Pensé qué hace especial a cada app antes de decidir.

\begin{infoband}{milcVerde}
\textbf{Lo que va al cuaderno (Actividad 1 · IDENTIFICA).} Título: «Actividad 1 --- 9 situaciones, 3 apps». \textbf{Formato:} lista numerada del 1 al 9 con su letra (W/E/P). \textbf{Extensión:} 9 líneas. \textbf{Sabes que terminaste cuando} cada situación tiene una sola letra asignada.
\end{infoband}

\puente{Después aprendes qué hace cada app y dónde están sus funciones esenciales.}

\begin{titlebox}{milcTurquesa}
Fase 2 · Sistematización con pensamiento computacional
\end{titlebox}

\begin{softbox}{milcTurquesa}{milcGris}{Cuatro pilares aplicados al tema}
Tu intuición se afila ahora con la sistematización. Aprende qué hace cada app y dónde están sus 5 funciones esenciales. Las respuestas correctas de la Actividad 1: \emph{1=W, 2=E, 3=P, 4=W, 5=E, 6=P (puede ser W con diseño), 7=W, 8=E, 9=P (con datos de Excel)}.
\end{softbox}

\small
\begin{tabularx}{\textwidth}{>{\bfseries\RaggedRight\arraybackslash}p{3.6cm}Y}
\rowcolor{milcTurquesa!90}\color{white}Pilar & \color{white}\bfseries Aplicación al tema\\
1. Descomposición & \textbf{Word Online --- el lápiz digital.} \textbf{Para qué sirve:} escribir textos. Ensayos, cartas, cuentos, informes, poemas, cualquier cosa con palabras. \textbf{5 funciones esenciales (ya las conoces de G6·P3):} (1) Formato de texto (negrita, alineación, color). (2) Listas con viñetas y numeradas. (3) Tablas. (4) Insertar imagen. (5) Encabezado y pie de página. \textbf{Cómo abrir:} \texttt{office.com} → ícono Word (W azul). Crea documento en blanco. \textbf{Guardado:} automático en OneDrive cada pocos segundos.\\
2. Reconocimiento de patrones & \textbf{Excel Online --- la calculadora maestra.} \textbf{Para qué sirve:} todo lo que tiene \emph{números organizados en tabla} (calificaciones, presupuestos, listas con conteos, estadísticas). \textbf{5 funciones esenciales:} (1) Celdas (cada cuadrito es una celda; se identifican A1, B2, C3...). (2) Fórmulas: empiezan con \texttt{=}, ejemplo \texttt{=A1+B1} suma. (3) Función \texttt{=PROMEDIO(rango)} calcula promedio. (4) Gráficos automáticos: seleccionas datos → Insertar → Gráfico. (5) Formato condicional: las celdas se pintan de color según el valor. \textbf{Cómo abrir:} \texttt{office.com} → Excel (X verde).\\
3. Abstracción & \textbf{PowerPoint Online --- el mostrador de ideas.} \textbf{Para qué sirve:} hacer presentaciones con diapositivas que se proyectan. Cada diapositiva es como una página visual: con título, texto corto, imágenes. \textbf{5 funciones esenciales:} (1) Nueva diapositiva (insertar más): se hace con un clic. (2) Plantillas (\emph{Diseño → Temas}): cambian colores y fuentes en bloque. (3) Insertar imagen, video, ícono. (4) Transiciones entre diapositivas (animaciones). (5) Modo presentador (F5): ves notas del orador mientras la audiencia ve solo las diapositivas. \textbf{Cómo abrir:} \texttt{office.com} → PowerPoint (P naranja).\\
4. Algoritmo conceptual & \textbf{Regla maestra para escoger app.} Hazte 3 preguntas. \textbf{(1) ¿La información son palabras corridas?} → Word. \textbf{(2) ¿La información son números o datos en tabla?} → Excel. \textbf{(3) ¿La información se va a presentar a otros visualmente?} → PowerPoint. Si la respuesta es \emph{``más de una''}, escoge la principal. Ejemplo: \emph{``encuesta con gráficos a presentar''} → datos en Excel, gráficos en PowerPoint, ambos colaborando.\\
\end{tabularx}
\normalsize

\begin{drawbox}{milcTurquesa}{3 apps + reglas de elección}
\textbf{Word:} para palabras (textos largos). \\[1mm]
\textbf{Excel:} para números (tablas y cálculos). \\[1mm]
\textbf{PowerPoint:} para presentar (diapositivas). \\[1mm]
\textbf{Regla maestra:} pregunta principal → ¿palabras, números o mostrar? → escoge la app.
\end{drawbox}

\begin{softbox}{milcMostaza}{milcGris}{Errores comunes y su corrección}
\textbf{Confusiones frecuentes al elegir.} (1) \emph{Hacer un horario en Word con tabla} --- válido pero limitado: Excel hace tabla con cálculos automáticos. (2) \emph{Hacer presentación en Word con saltos de página} --- pierde transiciones, animaciones, modo presentador. (3) \emph{Calcular promedio en Word a mano} --- Excel lo hace con \texttt{=PROMEDIO()}. (4) \emph{Escribir un cuento en PowerPoint} --- pierde el flujo narrativo del texto corrido. (5) \emph{Usar Excel para hacer una carta formal} --- las celdas distorsionan el formato de carta. Cada app para su uso.
\end{softbox}

\begin{infoband}{milcTurquesa}
\textbf{Lo que va al cuaderno (Actividad 2 · EXPLICA).} Título: «Actividad 2 --- Tabla comparativa Word/Excel/PowerPoint». \textbf{Formato:} tabla de 5 filas, 3 columnas (Word, Excel, PowerPoint). Filas: \emph{(1) ¿Para qué sirve?}, \emph{(2) Información típica}, \emph{(3) Función especial 1}, \emph{(4) Función especial 2}, \emph{(5) Ejemplo escolar}. \textbf{Extensión:} 5 filas. \textbf{Sabes que terminaste cuando} podrías recomendarle a un primo la app correcta para 5 tareas distintas.
\end{infoband}

\puente{Con todo claro, creas 1 Word + 1 Excel + 1 PowerPoint y los guardas correctamente.}

\begin{titlebox}{milcMagenta}
Fase 3 · Praxis con pensamiento computacional
\end{titlebox}

\begin{softbox}{milcMagenta}{milcGris}{Construyendo el producto con los cuatro pilares}
\textbf{Actividad 3 · CREA --- 1 Word + 1 Excel + 1 PowerPoint reales} (32 min · individual). Vas a crear los 3 archivos en línea y dejarlos guardados en tu OneDrive en la carpeta correcta. Si no tienes computador en clase, planificas el contenido en el cuaderno y los creas en casa.
\end{softbox}

\small
\begin{tabularx}{\textwidth}{>{\bfseries\RaggedRight\arraybackslash}p{3.6cm}Y}
\rowcolor{milcMagenta!90}\color{white}Pilar & \color{white}\bfseries Aplicación al producto\\
Descomposición & \textbf{Paso 1 --- Abre office.com.} En el navegador, inicia sesión con tu cuenta institucional. Verás la página principal de Microsoft 365 con todas las apps. Ahora vas a crear 3 archivos seguidos en esta sesión.\\
Patrones & \textbf{Paso 2 --- Word en línea (10 min).} Haz clic en el ícono Word (W azul). Crea \emph{Documento en blanco}. En la parte superior izquierda, donde dice \emph{``Documento1''}, haz clic y renómbralo: \texttt{2026-mi-vision-de-la-minga-digital}. Escribe \textbf{1 página} sobre cómo entiendes el trabajo colaborativo en la nube después de las primeras 3 sesiones. 4 párrafos cortos: \emph{(1) ¿qué es la minga?}, \emph{(2) ¿cómo se parece a Microsoft 365?}, \emph{(3) ¿qué he activado y aprendido?}, \emph{(4) ¿qué espero del periodo?}. Aplica formato (justificado, negrita en ideas importantes, sangría). Cuando termines, ve a la pestaña \emph{Mover a} y mueve el archivo a tu carpeta \texttt{Mis-documentos/Colegio/Tecnologia/Periodo-1}.\\
Abstracción & \textbf{Paso 3 --- Excel en línea (10 min).} Vuelve a office.com, ahora Excel (X verde). Crea \emph{Libro en blanco}. Renombra: \texttt{2026-calificaciones-prueba}. Crea una tabla con \textbf{5 columnas} (encabezados: Nombre, Sociales, Mate, Tecnología, Promedio) y \textbf{10 filas} de estudiantes inventados con notas inventadas (1-5). En la columna \emph{Promedio}, en la celda E2, escribe \texttt{=PROMEDIO(B2:D2)} y presiona Enter. Aparece el promedio de la primera fila. \emph{Truco mágico:} agarra el cuadrito de la esquina inferior derecha de E2 y arrástralo hacia abajo hasta E11: la fórmula se copia y calcula los promedios de todos. Muévelo a la misma carpeta.\\
Algoritmo de armado & \textbf{Paso 4 --- PowerPoint en línea (12 min).} office.com → PowerPoint (P naranja) → \emph{Presentación en blanco}. Renombra: \texttt{2026-mis-3-apps-favoritas}. Crea \textbf{5 diapositivas}: \emph{(D1) Portada}: título + tu nombre + grado. \emph{(D2) Word}: dile a la audiencia por qué te gusta Word. \emph{(D3) Excel}: lo mismo con Excel. \emph{(D4) PowerPoint}: lo mismo con PowerPoint. \emph{(D5) Cierre}: \emph{``¿Y tú cuál usarás más?''}. Aplica un tema bonito en \emph{Diseño}. Inserta al menos 1 imagen libre de Pixabay. Muévelo a la misma carpeta.\\
\end{tabularx}
\normalsize

\begin{drawbox}{milcMagenta}{Antes de cerrar la S4}
\checkbox Tengo 1 archivo Word en línea con texto de 1 página formateado. \\[1mm]
\checkbox Tengo 1 archivo Excel con tabla 5x10 + fórmula =PROMEDIO() aplicada. \\[1mm]
\checkbox Tengo 1 archivo PowerPoint con 5 diapositivas y diseño temático. \\[1mm]
\checkbox Los 3 archivos están en Mis-documentos/Colegio/Tecnologia/Periodo-1. \\[1mm]
\checkbox Cada archivo tiene nombre que sigue las 5 reglas. \\[1mm]
\checkbox Verifiqué en el cuaderno los 3 archivos creados.
\end{drawbox}

\begin{softbox}{milcMagenta}{milcGris}{Plantilla de trabajo}
\textbf{Plan de la S4.} \textit{Paso 1:} office.com. \textit{Paso 2:} Word 1 página sobre minga digital. \textit{Paso 3:} Excel calificaciones con PROMEDIO. \textit{Paso 4:} PowerPoint 5 diapositivas sobre 3 apps. \textit{Paso 5:} todos en Tecnologia/Periodo-1. \textit{Paso 6:} registro y firma en cuaderno.
\end{softbox}

\begin{infoband}{milcMagenta}
\textbf{Lo que va al cuaderno (Actividad 3 · CREA).} Título: «Actividad 3 --- Mis 3 archivos estrella». \textbf{Formato:} registro de los 3 archivos + reflexión sobre cuál te gustó más. \textbf{Extensión:} 1 página. \textbf{Sabes que terminaste cuando} tus 3 archivos están vivos en OneDrive y puedes mostrarlos al profe.
\end{infoband}

\puente{El producto son los 3 archivos en OneDrive + tabla comparativa en cuaderno.}

\begin{titlebox}{milcMostaza}
Producto de la sesión
\end{titlebox}
\textbf{3 archivos creados (Word + Excel + PowerPoint) en OneDrive · organizados}.
\begin{softbox}{milcMostaza}{milcGris}{Criterios de calidad}
(1) El Word tiene 1 página con 4 párrafos sobre la minga digital, formateado. (2) El Excel tiene tabla 5x10 con fórmula PROMEDIO aplicada y propagada. (3) El PowerPoint tiene 5 diapositivas con diseño temático + 1 imagen libre. (4) Los 3 archivos están en la carpeta correcta del OneDrive. (5) Los nombres siguen las 5 reglas (sin espacios, sin acentos, con fecha).
\end{softbox}

\puente{Antes de salir, verifica que los 3 archivos abren desde otro dispositivo (sincronización OK).}

\begin{titlebox}{milcVino}
Fase 4 · Evaluación liberadora — cinco dimensiones
\end{titlebox}

\begin{softbox}{milcVino}{milcGris}{Sentido de la evaluación}
Evaluar no es solo revisar una nota. Es reconocer qué aprendí, qué cambió en mí, qué emociones manejé, cómo me relaciono con la ciudad y la comunidad, y qué le dejaré a quien viene detrás.
\end{softbox}

\textbf{Mi reflexión liberadora en cinco dimensiones.} Completa cada frase iniciada:

\small
\begin{tabularx}{\textwidth}{>{\bfseries\RaggedRight\arraybackslash}p{3.4cm}Y>{\RaggedRight\arraybackslash}p{4.6cm}}
\rowcolor{milcVino}\color{white}Dimensión & \color{white}\bfseries Mi respuesta (frame starter) & \color{white}\bfseries Algo que descubrí\\
1. Personal & Lo que más cambió en mí fue \linead{6.5cm} & \linead{4.2cm}\\
2. Emocional & La emoción más fuerte fue \linead{2.5cm} y la regulé \linead{3cm} & \linead{4.2cm}\\
3. Ciudadana & En democracia esto importa porque \linead{6cm} & \linead{4.2cm}\\
4. Local (Cartago) & En mi barrio sería útil para \linead{6.5cm} & \linead{4.2cm}\\
5. Intergeneracional & A mi abuela/abuelo le diría \linead{6.5cm} & \linead{4.2cm}\\
\end{tabularx}
\normalsize

\begin{infoband}{milcVino}
\textbf{Para la próxima sustentación.} Vuelve a esta tabla en seis meses. Verás que las respuestas cambian — no porque lo aprendido cambie, sino porque tú cambias. La evaluación liberadora es una conversación contigo a través del tiempo.
\end{infoband}


\puente{Cerramos con 3 voces que te ayudan a entender por qué dominar pocas herramientas a fondo vale más que muchas a medias.}

\begin{titlebox}{milcGradeDark}
Triángulo de pensamiento · cierre filosófico
\end{titlebox}

\begin{softbox}{milcVino}{milcGris}{Dussel · lente del nosotros}
\textit{````Dominar 3 herramientas a fondo te hace artesano. Saber 10 a medias te hace turista.''''}

Hay personas que se jactan de \emph{``conocer todo Microsoft 365''} pero en realidad no dominan ninguna app a fondo. Don Lucho prefería que sus aprendices \textbf{dominaran 3 herramientas estrella} antes de querer aprender 10. Esa disciplina antigua aplica a Microsoft 365: \emph{Word, Excel, PowerPoint} bien aprendidos te dan más poder real que las 8 a medias.

\textbf{¿Quiero ser turista del software o artesano que domina lo esencial?}
\end{softbox}
\linead{\linewidth}\\[1mm]
\linead{\linewidth}

\begin{softbox}{milcMostaza}{milcGris}{Estoicismo · Marco Aurelio (emperador que entendía la economía de medios) · cuidado interior}
\textit{````Pocas armas bien usadas ganan más batallas que muchas mal manejadas.''''}

Marco Aurelio dirigía un ejército. Sabía que \textbf{soldados bien entrenados en pocas armas son más letales que mediocres con muchas}. Tu vida adulta es batalla constante (en sentido amplio): tareas, trabajos, decisiones. \emph{Dominar Word, Excel y PowerPoint a fondo te da ventaja sobre el 80\% de tus pares que las usan a medias}. Vale la pena la inversión inicial.

\textbf{¿Estoy aprendiendo lo esencial a fondo, o coleccionando conocimientos superficiales?}
\end{softbox}
\linead{\linewidth}\\[1mm]
\linead{\linewidth}

\begin{softbox}{milcTurquesa}{milcGris}{Floridi · lente de la infoesfera}
\textit{````El que aprende Word, Excel y PowerPoint en serio se ahorra años de improvisación profesional.''''}

En el trabajo y la universidad, \textbf{las 3 apps se exigen sin enseñar}: se asume que las sabes. Los que las aprendieron en colegio (como tú) avanzan sin pena. Los que llegan sin saberlas pasan meses sufriendo. \emph{Hoy aprendiste lo que el 60\% de adultos colombianos no sabe}. Eso es ventaja real. Aprovéchala.

\textbf{¿Veo este aprendizaje como un favor que me hace mi yo de 25 años?}
\end{softbox}
\linead{\linewidth}\\[1mm]
\linead{\linewidth}

\begin{titlebox}{milcVino}
Compromiso final
\end{titlebox}

\small
\begin{tabularx}{\textwidth}{>{\RaggedRight\arraybackslash}p{3.0cm}YYY}
\rowcolor{milcVino}\color{white}\bfseries Criterio & \color{white}\bfseries Lo logré & \color{white}\bfseries En proceso & \color{white}\bfseries Necesito apoyo\\
Comprensión del tema & Explico ideas centrales con mis palabras. & Reconozco ideas pero falta precisión. & Necesito repasar conceptos base.\\
Pensamiento computacional & Apliqué los 4 pilares al diseñar mi producto. & Apliqué algunos pilares. & Necesito releer la fase 2.\\
Aplicación y producto & Mi producto cumple los criterios y se puede revisar. & Producto funcional con ajustes pendientes. & Aún en construcción.\\
Reflexión 5 dimensiones & Respondí cada dimensión con honestidad. & Respondí algunas con detalle. & Mis respuestas son muy generales.\\
Triángulo de pensamiento & Conecté las 3 voces con mi trabajo real. & Conecté al menos una. & No relaciono las voces con mi proceso.\\
\end{tabularx}
\normalsize

\textbf{Hoy entendí que:}\\[1mm]
\linead{\linewidth}\\[1mm]
\linead{\linewidth}

\textbf{Mi compromiso para la próxima sesión es:}\\[1mm]
\linead{\linewidth}\\[1mm]
\linead{\linewidth}

\begin{softbox}{milcMostaza}{milcGris}{Cita de despedida · Marco Aurelio}
\textit{``El obstáculo en el camino se vuelve el camino. Lo que estorba al camino se vuelve camino.''}\\[1mm]
Cada error de hoy, bien observado, se convierte en pieza del camino que pisarás mañana.
\end{softbox}

\small
\begin{tabularx}{\textwidth}{>{\bfseries}p{3.6cm}Y}
\rowcolor{milcGradeDark!12}Nombre completo: & \linead{10cm}\\
Fecha: & \linead{4cm} \quad Curso/grupo: \linead{4cm}\\
Mi firma: & \linead{9cm}\\
\end{tabularx}
\normalsize

\begin{infoband}{milcGradeDark}
\textbf{Para la próxima sesión.} Esta semana, vuelve a abrir tus 3 archivos al menos 1 vez para revisarlos y mejorarlos. Familiarízate con la sensación de \emph{guardado automático} (no necesitas Ctrl+G). La próxima clase aprendes a \textbf{compartir tus archivos} con compañeros: el primer paso de la minga digital real.
\end{infoband}

\end{document}
